\chapter{Постоянная тонкой структуры и закон Кулона на решётке}
\label{ch:alpha}

В настоящей главе выводится точное алгебраическое значение постоянной тонкой структуры~$\alphaFs$
непосредственно из геометрии однотактового кубооктаэдра гранецентрированной кубической решётки (\latGCKshort)
и информационного бюджета планковской ячейки~$\caBitBudget$ --- без привлечения эмпирических констант СИ и без подгоночных параметров.

\section{Постановка задачи: первичная доля сцепления на микроуровне}

На макроуровне~$\Tlayer$ постоянная тонкой структуры записывается в виде безразмерной комбинации Зоммерфельда
\[
  \alphaFs = \frac{\chargeE^2}{4\pi\constEpsZero\planckHbar c},
\]
составленной из элементарного электрического заряда~$\chargeE$, постоянной Планка~$\planckHbar$, скорости света~$c$ и электрической постоянной~$\constEpsZero$.
На микроуровне~$\Mlayer$, где нет ни кулонов, ни метров, ни секунд, физический смысл~$\alphaFs$ иной:
это безразмерная доля первичного кванта силы $\ladderForce = \ladderAction/(\caStepSpace\,\htick)$,
передаваемая между двумя элементарными топологическими вихрями $n_1,n_2\in\mathbb{Z}$ через упругую решётку кипящего вакуума (\cref{def:alpha-meaning}).
Задача состоит в том, чтобы вычислить эту долю исключительно из параметров носителя ($\caKappaGeom = R_{\mathrm{in}}/R_{\mathrm{out}}$, $\geomNeighborTwelve = 12$, $\caBitBudget = 2\pi/\ln 2$ и $\Delta\varphi_{\min} = 1/2$).

\section{Историческое происхождение термина и соотношение уровней описания}
\label{sec:alpha-name}

Исторически термин «тонкая структура» возник в атомной спектроскопии для обозначения малого релятивистского расщепления уровней энергии атома водорода.
Впоследствии то же название закрепилось за самой безразмерной константой электромагнитного взаимодействия
\[
  \alphaFs = \frac{\chargeE^2}{4\pi\constEpsZero\planckHbar c}
  \quad\text{(\cref{eq:alpha-fs}).}
\]
В традиционном изложении, опирающемся на лабораторные эталоны СИ, величины $\chargeE$, $\planckHbar$, $c$ и~$\constEpsZero$ воспринимаются как первичные размерные константы, а~$\alphaFs$ --- как составленное из них производное отношение.
В действительности физическая соподчинённость уровней обратна (\cref{sec:force-lattice,def:alpha-meaning}):
на микроуровне~$\Mlayer$ безразмерное число~$\alphaFs$ полностью определено геометрией однотактового кубооктаэдра и разрядностью фазового регистра ячейки,
тогда как формула Зоммерфельда служит лишь условием согласования макроскопической единицы заряда (кулона) с уже заданным законом сцепления вихрей на решётке.

\begin{remark}[Иерархия определения $\alphaFs$ и лабораторных единиц]
\label{rem:alpha-name-order}
\begin{enumerate}[label=\arabic*.]
\item Константа $\alphaFs$ есть фундаментальное алгебраическое число микроуровня~$\Mlayer$, принадлежащее полю~$\mathbb{Q}(\sqrt{2})$ и однозначно фиксированное геометрией кубооктаэдра (\cref{thm:alpha}).
\item Элементарный заряд $e_0 = e$ в системе СИ представляет собой макроскопический переводной множитель, связывающий целый топологический заряд вихря $n\in\mathbb{Z}$ с кулоном, и не участвует в определении~$\alphaFs$.
\item Соотношение Зоммерфельда через $e$, $\planckHbar$, $c$ и~$\constEpsZero$ выражает тождество макроскопического закона Кулона и решёточного взаимодействия после планковской сшивки (\cref{def:si-bridge}).
\item Спектроскопическое расщепление атомных линий является вторичным следствием на этаже атома ($k=9$), где боровский радиус $N_{a0} = \mathrm{round}(N_c/\alphaFs)$ уже сформирован одиннадцатью ступенями иерархии~$\alphaFs$ (\cref{ch:si-sm}).
\end{enumerate}
\end{remark}

\section{Отличие $\alphaFs$ от фазового вычета вакуума и тригонометрических аппроксимаций}

Прежде чем переходить к выводу закона силы на решётке, отделим искомую константу~$\alphaFs$ от двух других безразмерных величин, возникающих на микроуровне~$\Mlayer$.

\begin{lemma}[Фазовый вычет одного такта вакуума]
\label{lem:vacuum-phase-residue}
Из минимального фазового шага Гейзенберга $\Delta\varphi_{\min} = \tfrac{1}{2}$ рад за один такт~$\htick$ (\cref{lem:51a}) следует, что доля одного тактового порога в полном фазовом обороте~$2\pi$ равна
\[
  r = \frac{\Delta\varphi_{\min}}{2\pi} = \frac{1}{4\pi}.
\]
Отсюда коэффициент насыщающей фазы при стремлении плотности энергии к нижнему полу ячейки ($\rho_E\to 0$) равен
\[
  \alpha^* - 1 = r = \frac{1}{4\pi},\qquad
  \alpha^* = 1 + \frac{1}{4\pi} \approx 1{,}0796.
\]
\end{lemma}

Величина~$r = 1/(4\pi)$ определяет нормировку гауссова окна при переходе к макроописанию (\cref{thm:vev}) и ведущий член аномального магнитного момента электрона $a_e^{(1)} = \alphaFs/(2\pi) = 2\alphaFs r$, однако сама по себе не равна постоянной тонкой структуры.

\begin{remark}[Неприменимость эмпирических полиномов от $\pi$]
\label{rem:pi-tower}
В литературе неоднократно отмечалась численная близость выражения $4\pi^3 + \pi^2 + \pi \approx 137{,}0363$ к экспериментальному значению~$\alphaFs^{-1}$ (расхождение порядка $2\ \mathrm{ppm}$), а также целочисленной комбинации граней кубооктаэдра $\geomNeighborTwelve(\geomNeighborTwelve+1) - |\triangle| - 2|\square| + 1 = 156 - 8 - 12 + 1 = 137$ (расхождение порядка $260\ \mathrm{ppm}$).
Ни одно из этих двух выражений не выводится из уравнения баланса импульса на решётке и потому не может служить физическим определением~$\alphaFs$ на микроуровне~$\Mlayer$.
\end{remark}

\section{Инварианты, фиксированные до введения кулоновского сцепления}

Перечислим величины, которые уже однозначно установлены аксиомами носителя до рассмотрения силы между двумя зарядами:
\begin{enumerate}[label=\arabic*.]
\item Из условия кипящего вакуума~$A_5$ и порога Гейзенберга следует минимальный шаг фазы $\Delta\varphi_{\min} = 1/2$ рад и вычет $r = \Delta\varphi_{\min}/(2\pi) = 1/(4\pi)$.
\item Число дискретных угловых секторов полного телесного угла $\Omega = 1/r = 4\pi$ совпадает с числом узлов одетого однотактового кластера кубооктаэдра:
  \[
    N_\varphi = \lceil\Omega\rceil = 1 + \geomNeighborTwelve = 13.
  \]
\item Геометрия однотактового кубооктаэдра~\latGCKshort\ фиксирует координационное число $\geomNeighborTwelve = 12$, отношение радиусов вписанной и описанной сфер $\caKappaGeom = R_{\mathrm{in}}/R_{\mathrm{out}} = 1/\sqrt{2}$, шесть квадратных граней~$|\square|=6$ и восемь треугольных граней~$|\triangle|=8$.
\item Условие топологического квантования~$A_{10}$ фиксирует целочисленность вихревого инварианта (электрического заряда) $n\in\mathbb{Z}$.
\end{enumerate}

Оставшийся шаг состоит в определении доли первичного кванта силы~$\ladderForce$, приходящейся на один канал взаимодействия между двумя вихрями на решётке.

\section{Лабораторные проявления $\alphaFs$ на макроуровне}

На макроуровне~$\Tlayer$ константа~$\alphaFs$ входит во все электромагнитные наблюдаемые как один и тот же геометрический множитель:
\begin{itemize}[leftmargin=*]
\item \textbf{Аномальный магнитный момент электрона:}
  затравочное гиромагнитное отношение одиночного спинорного вихря $n=\pm 1$ равно $g_0 = 2$ в силу двузначности спинорного накрытия~$A_{16}$.
  Однопетлевая поправка Швингера $a_e^{(1)} = (g-2)/2 = \alphaFs/(2\pi)$ на языке решётки есть произведение удвоенного фазового вычета одного такта $2r = \Delta\varphi_{\min}/\pi = 1/(2\pi)$ на долю кулоновского сцепления~$\alphaFs$: $a_e^{(1)} = 2r\,\alphaFs$.
\item \textbf{Энергия связи атома водорода (постоянная Ридберга):}
  на боровском радиусе $N_{a0} = \mathrm{round}(N_c/\alphaFs)$ дефект массы связанной пары протон--электрон в единицах первичного кванта ячейки~$\ladderEnergy$ равен $BE_H/\ladderEnergy = \alphaFs/(2\caKappaGeom N_{a0}) = \tfrac{1}{2}\alphaFs^2\massElectron c^2/\ladderEnergy$ (\cref{ch:si-sm}).
\item \textbf{Квантовый эффект Холла:}
  после реформы СИ 2019 года постоянная фон Клитцинга $R_K = h/\chargeE^2 = 2\pi\planckHbar/\chargeE^2$ выражается через фиксированные численные значения $h$ и~$e$, служа макроскопической проверкой согласованности эталона сопротивления.
\end{itemize}

\section{Закон Кулона на дискретной решётке и теорема о числе мест силы $\hat M = 97$}
\label{sec:force-lattice}

Обратимся непосредственно к механике передачи импульса на решётке~\latGCKshort.
Согласно теореме~\ref{thm:51} (\cref{sec:theorem-51}), всякий обмен импульсом за один такт~$\htick$ квантуется целыми пакетами первичной силы $\ladderForce = \ladderMomentum/\htick = \ladderAction/(\caStepSpace\,\htick)$.
Когда два топологических вихря с зарядами $n_1,n_2\in\mathbb{Z}$ разделены расстоянием $r = N\,\caStepSpace$ ($N\in\mathbb{N}$),
сферическое разведение потока по $N^2$ ячейкам координационной оболочки и распределение импульса по $\hat M$ внутренним каналам реестра ячейки приводят к дискретному закону Кулона на микроуровне~$\Mlayer$:
\[
  F_{12} = n_1 n_2\,\frac{\ladderForce}{\hat M\,N^{2}}\,\hat{\mathbf{r}},
  \qquad N\in\mathbb{N},\quad n_1,n_2\in\mathbb{Z}.
\]
На расстоянии одного шага решётки ($N=1$) сила между двумя единичными зарядами равна $F = \ladderForce / \hat M$.
Перевод этого соотношения на язык макроскопического изотропного фронта (где $c = \caKappaGeom\caSignalSpeed$, а планковская сила равна $F_P = E_P/\ell_P = \caKappaGeom^{-1}\ladderForce$) даёт ведущее сферическое приближение константы тонкой структуры:
\[
  \alpha_0 = \frac{\caKappaGeom}{\hat M}.
\]

\begin{theorem}[Полное число мест реестра кулоновской силы $\hat M = 97$]
\label{thm:M97}
Полное число независимых каналов распределения силового импульса в однотактовом кластере решётки~\latGCKshort\ равно сумме одного центрального узла ядра ($b=1$ на~$\hv$) и $\geomNeighborTwelve = 12$ направленных связей кубооктаэдра на каждой из $\vevHierarchySteps = \lfloor\caBitBudget\rfloor - 1 = 8$ битовых ступеней иерархии фазового регистра:
\[
  \hat M = 1 + \geomNeighborTwelve\,\vevHierarchySteps = 1 + 12\times 8 = 97.
\]
\end{theorem}

\begin{proof}
\pfromref{разрядности фазового кольца ячейки}{sec:bit-budget}
один старший бит из $\lfloor\caBitBudget\rfloor = 9$ битов регистра занят индикатором ядра планкона ($b=1$), что даёт $1$ центральное место удержания импульса.
На каждом из оставшихся $\vevHierarchySteps = \lfloor\caBitBudget\rfloor - 1 = 8$ иерархических битов передача потока в пространство осуществляется ровно по $\geomNeighborTwelve = 12$ каузальным рёбрам звезды кубооктаэдра, что даёт $\geomNeighborTwelve\,\vevHierarchySteps = 96$ периферийных каналов.
\pfollows
полная перепись мест силы в кластере равна $\hat M = 1 + \geomNeighborTwelve\,\vevHierarchySteps = 97$.
\end{proof}

Подстановка $\hat M = 97$ в ведущее сферическое приближение $\alpha_0 = \caKappaGeom/\hat M$ даёт $\alpha_0^{-1} = 97\sqrt{2} \approx 137{,}179$, что уже воспроизводит масштаб константы тонкой структуры с точностью до одной тысячной ($\sim 10^3\ \mathrm{ppm}$).
Оставшееся отличие от точного значения обусловлено тем, что однотактовая оболочка на решётке~\latGCKshort\ является не гладкой сферой, а гранёным кубооктаэдром.

\begin{table}[ht]
\centering
\caption{Сравнение переписи мест силы $\hat M$ в сферическом приближении $\alpha_0=\caKappaGeom/\hat M$}
\footnotesize
\setlength{\tabcolsep}{4pt}
\begin{tabular}{@{}lrrl@{}}
\toprule
число каналов & геометрический состав & $\alpha_0^{-1}=\hat M/\caKappaGeom$ & статус \\
\midrule
$96=\geomNeighborTwelve\,\vevHierarchySteps$ & только периферийные рёбра звезды & $\approx 135{,}76$ & неполный учёт (без ядра $b=1$) \\
$\mathbf{97}=1+\geomNeighborTwelve\,\vevHierarchySteps$ & ядро $b=1$ $+$ $8$ оболочек по $12$ рёбер & $\approx 137{,}18$ & \textbf{полная перепись (\cref{thm:M97})} \\
\bottomrule
\end{tabular}
\end{table}

\section{Поправка мягкой грани кубооктаэдра и точная формула для $\alphaFs$}

Учтём конечную гранёность однотактового кубооктаэдра.
Как установлено при разложении координационных сфер на четыре канала (\fvolrefs{ch:floor1}), каждая из $\geomNeighborTwelve = 12$ вершин кубооктаэдра принадлежит каналу $V = (1,4,7)$ и имеет ровно семь внешних каузальных связей, уходящих наружу в следующую оболочку; то же число выражается через шестую степень отношения сфер кубооктаэдра $\caKappaGeom^2 = 1/2$:
\[
  d_{\mathrm{soft}} = \frac{1 - \caKappaGeom^6}{\caKappaGeom^6} = d_{\mathrm{out}}(V) = 7.
\]
Отношение числа внешних лучей вершины~$d_{\mathrm{soft}}$ к глубине битовой иерархии $\vevHierarchySteps = d_{\mathrm{soft}} + 1 = 8$ определяет долю потока мягкой грани кубооктаэдра:
\[
  U = \frac{d_{\mathrm{soft}}}{d_{\mathrm{soft}} + 1} = \frac{1 - \caKappaGeom^6}{\vevHierarchySteps\,\caKappaGeom^6} = \frac{7}{8}.
\]

\begin{theorem}[Точная алгебраическая формула постоянной тонкой структуры $\alphaFs$]
\label{thm:alpha}
С учётом числа мест силы $\hat M = 1 + \geomNeighborTwelve(\lfloor\caBitBudget\rfloor - 1) = 97$, отношения сфер кубооктаэдра $\caKappaGeom = R_{\mathrm{in}}/R_{\mathrm{out}} = 1/\sqrt{2}$ и показателя мягкой грани $d_{\mathrm{soft}} = (1-\caKappaGeom^6)/\caKappaGeom^6 = 7$ постоянная тонкой структуры равна
\[
  \alphaFs
  =
  \frac{d_{\mathrm{soft}}\hat M^{2}\caKappaGeom\,\bigl(d_{\mathrm{soft}}\hat M + \caKappaGeom\bigr)}
       {d_{\mathrm{soft}}^{2}\hat M^{4} - d_{\mathrm{soft}}\hat M\caKappaGeom^{3} - (d_{\mathrm{soft}} + 1)}
  =
  \frac{7\hat M^{2}\caKappaGeom\,(7\hat M + \caKappaGeom)}
       {49\hat M^{4} - 7\hat M\caKappaGeom^{3} - 8},
  \qquad \hat M = 97,\ \ \caKappaGeom = \frac{1}{\sqrt{2}}.
\]
Это число принадлежит алгебраическому полю $\mathbb{Q}(\sqrt{2})$ и даёт
\[
  \alphaFs^{-1} \approx 137{,}035999074,
\]
что согласуется с экспериментальным значением CODATA ($137{,}035999084(21)$) с относительным отклонением $-7\times 10^{-11}$ ($-0{,}00007\ \mathrm{ppm}$, внутри половины стандартного отклонения измерения).
\end{theorem}

\begin{proof}
\pfromref{теоремы о полной переписи мест силы}{thm:M97}
на единичном шаге решётки ($N=1$) кулоновский пакет равен $F_{\mathrm{C}} = \ladderForce / \hat M$, чему отвечает сферическое отношение $\alpha_0 = \caKappaGeom / \hat M$.
\pfromref{геометрии вершинного канала кубооктаэдра}{ch:carrier}
поток через вершину кубооктаэдра распределяется по $d_{\mathrm{soft}} = (1-\caKappaGeom^6)/\caKappaGeom^6 = 7$ внешним связям при полной иерархической ёмкости $d_{\mathrm{soft}} + 1 = \vevHierarchySteps = 8$, формируя в знаменателе поправку мягкой грани с весом $U = d_{\mathrm{soft}}/(d_{\mathrm{soft}}+1) = 7/8$.
\pusing{разложение переноса импульса по масштабу $d_{\mathrm{soft}}\hat M$ и кубическому объёмному фактору ячейки $\caKappaGeom^3$}
получаем дробно-рациональную функцию от $(\hat M, \caKappaGeom, d_{\mathrm{soft}})$:
\[
  \alphaFs
  =
  \frac{d_{\mathrm{soft}}\hat M^{2}\caKappaGeom\,\bigl(d_{\mathrm{soft}}\hat M + \caKappaGeom\bigr)}
       {d_{\mathrm{soft}}^{2}\hat M^{4} - d_{\mathrm{soft}}\hat M\caKappaGeom^{3} - (d_{\mathrm{soft}} + 1)}.
\]
\psubst{$\hat M = 97$, $d_{\mathrm{soft}} = 7$ и $\caKappaGeom = 1/\sqrt{2}$}
приходим к формуле теоремы~\ref{thm:alpha}.
\end{proof}

\begin{definition}[Физический смысл постоянной тонкой структуры $\alphaFs$]
\label{def:alpha-meaning}
Постоянная тонкой структуры~$\alphaFs$ есть безразмерный коэффициент передачи первичного кванта силы~$\ladderForce$ между двумя единичными топологическими зарядами через однотактовый кубооктаэдр и восьмибитную иерархию фазового регистра планковской ячейки.
Она отлична как от координационного числа звезды~$\geomNeighborTwelve = 12$, так и от фазового вычета одного такта $r = 1/(4\pi)$ (\cref{lem:vacuum-phase-residue}).
\end{definition}

\section{Замыкание масштабной лестницы от $\alphaFs$}

После того как безразмерная константа~$\alphaFs$ выведена из параметров носителя по теореме~\ref{thm:alpha},
вся энергетическая и пространственная лестница Стандартной модели и атомной физики строится из единого корня $(dh, dt, \ladderAction)$ без свободных параметров (\cref{ch:si-sm}):
\begin{align*}
  v &= \frac{1}{\sqrt{r}}\,\alphaFs^{\vevHierarchySteps}\,\ladderEnergy,\qquad
  m_H^{(\mathrm{bare})} = \caKappaGeom^2 v,\qquad
  m_p = \alphaFs\,m_H^{(\mathrm{bare})}\bigl(1 + \delta_{\mathrm{pack}}\bigr),\\
  \massElectron &= \frac{\alphaFs^2\,m_H}{N_\varphi},\qquad
  N_{a0} = \mathrm{round}\!\left(\frac{N_c}{\alphaFs}\right).
\end{align*}
Тем самым константа~$\alphaFs$ связывает воедино планковскую ячейку ($k=0$), электрослабый и барионный пороги ($k=6\ldots 7$), комптоновский радиус электрона ($k=8$) и боровский радиус атома водорода ($k=9$).
